\section{A narrower arithmetic principle: successor rigidity}
\label{sec:successor-rigidity}

This section gives a proved reduction of an arithmetic principle, not
a second completed axiomatization.  The existing
\texttt{scope\_kernel} note obtains induction from scalar iterator
uniqueness by a retraction argument.  The incoming
\texttt{ipra\_scalar\_uniqueness} report completes the equality-reflection
bootstrap for that unrestricted uniqueness rule.  The observation here
is narrower: uniqueness just for zero-starting successor trajectories
already gives induction for every literal equation with zero on one side.
The proof uses one trajectory and does not require an idempotence lemma.

\subsection{Exact background and the restricted rule}

Use ordinary intuitionistic first-order logic with actual equality,
the scalar terms $0$, $S$, and $I(n,a;u.h)$, and the equations
\[
 I(0,a;u.h)=a,\qquad
 I(Sn,a;u.h)=h[I(n,a;u.h)/u].
\]
The bound variable $u$ occurs only in the body $h$.  Retain numerical
coverage
\[
                 \forall n\,(n=0\lor\exists m\,n=Sm).
\]
For the theorem below, neither successor injectivity nor separation of
$0$ from $S0$ is needed.  Define the selector
\[
 C(v,a,b):=I(v,a;u.b),\qquad
 C(0,a,b)=a,\quad C(Sv,a,b)=b.
\]

The proposed restricted arithmetic rule is
\[
 \frac{\Gamma\vdash f(0)=0\qquad
       \Gamma\vdash f(Sz)=S(f(z))}
      {\Gamma\vdash f(t)=t}.
 \tag{SR$_0$}\label{rule:sr-zero}
\]
Here $f(z)$ is one fixed finite scalar term template; it is not a
function variable.  The variable $z$ is fresh for $\Gamma$, $t$, and
all passive parameters.  In particular, the recurrence premise has
no local induction hypothesis.  This rule is narrower than arbitrary
scalar-iterator uniqueness: its seed is zero and its update is always
successor.  It does have a separate base premise, which must be counted
in any implementation comparison.

Literal-equation induction derives \eqref{rule:sr-zero}: induct on
$f(n)=n$.  The given base supplies the zero case.  Under $f(z)=z$,
the recurrence gives
$f(Sz)=S(f(z))=Sz$.  Unrestricted scalar uniqueness also derives the
rule: apply it to $f$ with update $S$ and to the identity template
$n\mapsto n$, then use the common zero seed.

\subsection{Induction by a trajectory that stays inside the predicate}

The following sufficient condition makes the proof transparent.  Let
$E(x)$ be a fixed formula and suppose that the retained context proves
its base and its successor step:
\[
 \Gamma\vdash E(0),\qquad
 \Gamma,E(z)\vdash E(Sz),
\]
where $z$ is fresh.  Suppose also that there is a fixed scalar term
$e(x)$ such that, uniformly in a fresh $x$,
\[
 \Gamma\vdash E(e(x)),\qquad
 \Gamma,E(x)\vdash e(x)=x.                         \tag{Reset}
\]
No decidability of the whole formula $E$ is included in this condition.

Define the finite term
\[
                     r(n):=I(n,0;u.e(Su)).
\]
We first prove $\Gamma\vdash E(r(n))$ using coverage on $n$.
If $n=0$, computation gives $r(n)=0$, and the base gives the result.
If $n=Sm$, computation gives $r(n)=e(S(r(m)))$, and the first
\textnormal{(Reset)} statement, instantiated at $S(r(m))$, gives the
result directly.  This argument does not assume $E(r(m))$ in its
successor branch and does not use numerical induction.

Now instantiate the given successor step at $r(z)$ and use the
already proved $E(r(z))$.  It gives $E(S(r(z)))$.  The second
\textnormal{(Reset)} statement consequently yields
\[
 e(S(r(z)))=S(r(z)).
\]
Thus the iterator equations prove, under the unchanged context,
\[
                   r(0)=0,\qquad r(Sz)=S(r(z)).
\]
One application of SR$_0$ gives $r(t)=t$.  Transport the theorem
$E(r(t))$ along this equation to obtain $E(t)$.

All substitutions in this proof are capture avoiding.  The eigenvariable
of SR$_0$ is chosen fresh for the retained context and passive parameters;
the auxiliary witness used when eliminating coverage is also fresh.
Consequently the argument works with arbitrary first-order assumptions
in $\Gamma$, including quantified assumptions.

\subsection{Unconditional consequence for zero equations}

Take $E(x)$ to be the literal equation $q(x)=0$ and assume its base
and successor step as above.  Put
\[
                       e(x):=C(q(x),x,0).
\]
If $q(x)=0$, selector computation gives $e(x)=x$, which proves the
second \textnormal{(Reset)} statement.  To prove the first, apply
coverage to $q(x)$.

If $q(x)=0$, then $e(x)=x$ and hence $q(e(x))=0$.  If
$q(x)=Sv$, then $e(x)=0$ and the assumed base $q(0)=0$ again gives
$q(e(x))=0$.  These are constructive numerical cases.  They require
neither a test for equality of two arbitrary terms nor separation of
zero and successor.

It follows that the stated background plus SR$_0$ derives the complete
contextual rule
\[
 \frac{\Gamma\vdash q(0)=0\qquad
       \Gamma,q(z)=0\vdash q(Sz)=0}
      {\Gamma\vdash q(t)=0}.
 \tag{ZI}\label{rule:rigidity-zero-induction}
\]
The construction of $e$ and $r$ is effective on finite syntax, so the
argument is a uniform finite proof transformation, not only a semantic
observation about the standard natural numbers.

\subsection{The exact unresolved step}

This result does not establish full scalar iPRA from SR$_0$.
The missing bootstrap is substantial.  A term $\delta(a,b)$ satisfying
\[
                 \delta(a,b)=0\quad\longleftrightarrow\quad a=b
\]
throughout the numerical domain would transfer ZI to induction on
every literal equation.  The already established scalar bridge would
then give the desired bi-interpretation with iPRA.  However, the current
predecessor and equality-reflection proofs use inductions on general
equations, or uniqueness with updates other than successor.  Importing
them before establishing this bridge would be circular.

For example, the completed unrestricted-uniqueness predecessor proof
compares unbounded trajectories with updates such as $v\mapsto v+3$,
$v\mapsto v+M(v)+4$, and its factor map $H$.  Those are not instances
of SR$_0$.  Boolean finite-state invariants can often be translated to
ZI, but that alone does not justify those unbounded trajectory
equalities.  No converse ZI$\Rightarrow$SR$_0$, no full equivalence with
unrestricted uniqueness, and no smaller complete checker for SR$_0$
are claimed here.

The precise further research question is whether scalar computation,
coverage, separation, and SR$_0$ construct a faithful equality test
without using any unrestricted equality induction.  A positive answer
would remove the arbitrary update template from the arithmetic
uniqueness rule; its two-premise implementation would then need a
matched size comparison.
